
\usepackage{lscape}
\usepackage{amsmath, amssymb, amsthm,amsfonts}
\usepackage{times}  %Usa a fonte Latin Modern
\usepackage[utf8]{inputenc} %Codificacao do documento
\usepackage[T1]{fontenc}    %Selecao de codigos de fonte.
\usepackage{indentfirst}%   Indenta o primeiro parágrafo de cada seção.
\usepackage{nomencl}    %Lista de simbolos
\usepackage{color}  %Controle das cores
\usepackage[table]{xcolor}  % colorir células numa tabela
\usepackage{multicol}   % Multiplas colunas
\usepackage{graphicx,subfig,float,wrapfig}    %Inclusão de gráficos e subfiguras 1a, 1b e 1c...
\usepackage{microtype}  %Para melhorias de justificação
\usepackage{lipsum} %Para geração de dummy text
\usepackage[nogroupskip,nonumberlist]{glossaries}   % Permite fazer o glossario
\usepackage[abnt-emphasize=bf,abnt-and-type=e,alf]{abntex2cite} % Citações ABNT
\usepackage{mathptmx}
\usepackage{appendix} % Gerar apêndice
\usepackage[bottom=2cm,top=3cm,left=3cm,right=2cm]{geometry}
% Configuracoes do documento

\usepackage{float} % Adicione o pacote float
\graphicspath{{./figuras/}} %Images na pasta "Figuras"

\setsecheadstyle{\bfseries \normalsize \uppercase}
\setsubsecheadstyle{\normalsize \uppercase}
\setsubsubsecheadstyle{\bfseries \normalsize}
\setlrmarginsandblock{3cm}{2cm}{*}  %Margens esquerda-direita
\setulmarginsandblock{3cm}{2cm}{*}  %Margens cima-baixo
\checkandfixthelayout
\setlength{\parindent}{1.25cm}  %paragrafo
\OnehalfSpacing %espacamento de 1,5
\setlength{\ABNTEXcitacaorecuo}{4cm}    %recuo citacao direta +3
\selectlanguage{brazil} % Seleciona o idioma do documento
\frenchspacing          % Retira espaço extra obsoleto entre as frases.
\numberwithin{section}{chapter} % Faz as seções serem numeradas com base no capítulo
\usepackage{orcidlink}
\usepackage{listings}	
\usepackage{tikz}

\usetikzlibrary{petri, positioning}
\usepackage[ruled,vlined]{algorithm2e}

\hypersetup{%
	colorlinks = true,
	linkcolor = blue, % Define a cor do texto
	citecolor = blue!85, % Define a cor da citação
}

\newacronym{RPH}{RPH}{Redes de Petri Hierárquicas}
\newacronym{RP-Tz}{RP-Tz}{Rede de Petri Temporizada}
\newacronym{RP-T}{RP-T}{Rede de Petri Temporal}
\newacronym{RPC}{RPC}{Redes de Petri Coloridas}
\newacronym{SED}{SED}{Sistemas a Eventos Discretos}
\newacronym{RP}{RP}{Redes de Petri}
\newacronym{IoT}{IoT}{\textit{Internet of Things} (Internet das coisas)}
\newacronym{AI}{AI}{\textit{Artificial Intelligence} (Inteligência Artificial)}
\newacronym{CAM}{CAM}{\textit{CAM - Computer Aided Manufacturing})  (Manufatura Assistida por Computador)}

\usepackage{enumitem} % Para o gantt graph
\usepackage{pgfgantt}

% \newglossaryentry{stakeholder}{
%      name=Stakeholder,
%      description = {qualquer indivíduo ou organização que, de alguma modo, é impactado pelas ações de uma determinada empresa. Em uma tradução livre para o português, o termo significa parte interessada.}
% }

\makeglossaries

\usepackage{xcolor}
\definecolor{codebg}{rgb}{0.95,0.95,0.95}
\lstset{
	backgroundcolor=\color{codebg},
	basicstyle=\ttfamily\footnotesize,
	breaklines=true,
	keywordstyle=\color{blue},
	stringstyle=\color{black},
	commentstyle=\color{orange},
	language=Python,
	showspaces=false,
	showstringspaces=false,
	columns=fullflexible,
	frame=single
}
% Comando para fonte das figuras e espaçamento

